% Metodologia. Fonte das decisões: Lab-03/docs/protocolo.md e config/estudo.json.
\section{Metodologia}
\label{sec:metodologia}

\subsection{População, seleção e amostra analisada}
\label{sec:amostra}

A população-alvo são repositórios públicos do GitHub com mais de 1.000
estrelas que usam GitHub Actions. A busca \texttt{stars:>1000}, ordenada por
estrelas, informou 65.512 repositórios. Os 1.800 mais populares formaram a
lista de candidatos. Destes, 1.563 usam Actions. São elegíveis os que têm, na
janela, pelo menos cinco releases publicadas no histórico do branch padrão e
pelo menos 50 workflow runs válidos disparados por \texttt{push} nesse branch.
Para limitar o custo da coleta, também ficaram de fora repositórios com mais de
500 releases ou mais de 1.000 runs de \texttt{push} na janela. Restaram 311
elegíveis, e os 300 mais populares formaram a seleção. A amostra é intencional
(por popularidade), não aleatória.

A coleta foi dividida nos índices 0--182 e 183--299. Na segunda parte, 17
repositórios não cumpriram o mínimo de releases na verificação de ancestralidade
e foram substituídos pelos próximos elegíveis em ordem de estrelas
(\texttt{data/selecao-amostra-revisada.json}).

\textbf{Amostra analisada.} O dataset congelado
(\texttt{data/amostra\_completa.json.gz}, coleta de 13 de outubro de 2026)
reúne 183 e 117 repositórios, totalizando os 300 selecionados. A janela e os
critérios são comuns às partes; todos usam a heurística v2. A amostra varia de
cerca de 24 mil a 289 mil estrelas e não representa projetos menos populares.

\subsection{Janela e unidades}

A janela vai de 6 de outubro de 2025 a 7 de outubro de 2026 (UTC$-3$), com fim
exclusivo, e tem 366 dias. O dia além dos doze meses é uma exceção de
compatibilidade, discutida na introdução. Releases e runs entram pela data de
publicação e de criação, respectivamente. Commits incluídos em releases da
janela podem ser anteriores a ela.

\subsection{Métricas}

\textbf{Frequência (RQ01):} releases publicadas na janela divididas pela sua
duração em semanas. \textbf{Lead time (RQ02):} cada release é comparada com a
anterior. A variante (a) é a mediana, por repositório, do intervalo entre a
publicação e o commit mais antigo da release. A variante (b) é a mediana sobre
todos os commits incluídos. \textbf{CFR (RQ03):} a variante (a) é a proporção de
runs de CI com falha entre os runs com sucesso ou falha. A variante (b) é a
proporção de releases seguidas, em até sete dias, por uma release corretiva.
Releases dos últimos sete dias da janela são censuradas. A heurística v2
considera corretiva a release que só altera o \textit{patch} semver, ou em que
pelo menos metade dos commits cita \texttt{revert}, \texttt{hotfix} ou
\texttt{fix}. \textbf{Recuperação (RQ04):} horas entre o início da primeira
falha e o próximo sucesso no mesmo workflow. Episódios abertos no fim da
janela são censurados.

A classificação DORA usa os cortes do enunciado. A categoria geral é a mediana
das quatro notas (Elite~4 a Low~1), arredondada para baixo, e só é emitida
quando as quatro métricas são observadas. A configuração de referência C1 usa
releases publicadas, lead time (a) e CFR (a). As configurações C2 (com
pré-releases, lead time (b) e CFR (b)) e C3 (tags) servem à RQ07. Ausências
nunca são imputadas como zero, e cada uma tem o motivo registrado.

\subsection{Validação manual}

O protocolo previa 60 repositórios, cinco releases por projeto e rotulagem
cega por três avaliadores. O relatório disponível cobre 283 releases em 58
repositórios, mas seu hash de origem difere do dataset congelado; os resultados
não validam diretamente o CFR (b) dos 300 repositórios. Só um avaliador humano
concluiu a rotulagem, sem consenso (Seção~\ref{sec:validacao-cfrb}).

\subsection{Análise estatística}

As métricas são descritas por mediana e IQR. A RQ05 usa a correlação de
Spearman entre a frequência e cada CFR. A RQ06 usa Kruskal-Wallis, com
$\varepsilon^2$ como tamanho de efeito, para linguagem principal e quartis de
estrelas e de contribuidores, fatores fixados antes da coleta. Comparações
múltiplas são corrigidas por Holm ($\alpha = 0{,}05$). A RQ07 compara as
classificações por percentual de mudança de categoria e kappa de Cohen
ponderado linear. Todas as análises são reexecutáveis a partir do dataset
congelado (\texttt{python -m lab03.analysis.*}).
